\BourbakiBackMatter{历史注记}

我们已经看到（第一章与第二、三章的历史注记），最初的非交换代数出现于 1843--44 年，见于 Hamilton {[}25{]} 与 Grassmann {[}23{]}、{[}24{]} 的著作中。Hamilton 在引入四元数时，就已经对实数域上的任意有限秩代数有了清晰的想法（{[}25{]}，序言，第 26--31 页）\(^{(1)}\)。在发展他的理论的过程中，他后来产生了考虑他所谓的“双四元数”的想法，即复数域上与四元数代数有相同乘法表的代数；在这个场合，他注意到这种扩张导致零因子的出现（{[}25{]}，第 650 页）。Grassmann 的观点略有不同，在很长一段时间里，他的“外代数”一直与一般代数理论相互分离；\(^{(2)}\) 但在他那尚欠精确的语言中，人们不能不辨认出由生成元组与关系定义的（实数域上的、有限维或非有限维的）代数的最初思想 {[}24{]}。

新的代数例子在 1850 年至 1860 年间被或多或少明确地引入。Cayley 在发展矩阵理论时 {[}8{]}，尚未把方阵视为构成一个代数（这一观点最先由 Peirce 父子于 1870 年前后清楚地表达 {[}40{]}）。不过，他已在这个场合注意到存在一个满足四元数乘法表的 2 阶矩阵组；这一观察可以看作代数的线性表示的第一个例子。\(^{(3)}\) 此外，在他定义有限群的抽象概念的那篇论文中，他还顺便给出了这样一个群的代数的定义，但没有进一步应用这个定义（{[}7{]}，第 129 页）。

1870 年之前没有别的进展，但正是在这个时候，对（实数域或复数域上的）有限维代数一般结构的研究开始了。B. Peirce 在这个方向上迈出了最初的步伐；他引入了幂零元素与幂等元素的概念，证明了一个（有或没有单位元的）代数只要含有一个非幂零元素就有一个非零幂等元，写下了著名的分解

\[
x = e x e + (x e - e x e) + (e x - e x e) + (x - x e - e x + e x e)
\]

（e 为一个幂等元，x 为任意元素），并且有了把一个幂等元分解为两两正交的“本原”幂等元之和的（仍略欠精确的）想法 {[}40{]}。此外，据 Clifford 所说（{[}11{]}，第 274 页）\(^{(4)}\)，两个代数的张量积这一概念应归功于 B. Peirce；Clifford 本人在若干年后把它隐式地应用于 Hamilton“双四元数”的一个推广 {[}9{]}，并显式地应用于以他的名字命名的代数的研究（{[}10{]} 与 {[}11{]}）。这些新概念被 B. Peirce 用来对（复数域上的）低维代数进行分类。这个问题在大约 1880 年也被英美学派的其他数学家处理过，特别是 Cayley 与 Sylvester。可能的结构五花八门这一事实很快显露出来，无疑正是这一事实使随后时期的研究转向获取具有更特殊性质的代数类。

在欧洲大陆，思想演进的方式颇为不同，这类研究在 1880 年之前即已出现。1878 年，Frobenius 证明四元数构成实数域上唯一的非交换（有限维）体的例子（{[}18{]}，第 59--63 页），这一结果两年后由 C. S. Peirce 独立发表 {[}39{]}。早在 1861 年，Weierstrass 在澄清 Gauss 的一个评注时，就在他的讲课中把 R 或 C 上不含幂零元素的交换代数 \(^{(5)}\) 刻画为同构于 R 或 C 的体的直和。Dedekind 在他的研究中，于 1870 年前后结合其交换域理论的“超复”观念得到了同样的结论。他们的证明发表于 1884--85 年（{[}54{]} 与 {[}12{]}）。也是在 1884 年，H. Poincaré 在一篇十分隐晦的注记 {[}41{]} 中，提请注意这样一种可能性：把方程组 \(z_{i} = \varphi_{i}(x_{1}, \ldots, x_{n}, y_{1}, \ldots, y_{n})\)（它表达一个代数中的乘法法则 \((\sum_{i} x_{i} e_{i})(\sum_{i} y_{i} e_{i}) = \sum_{i} z_{i} e_{i}\)）视为（当然是在局部）定义一个李群。这个评注似乎给 Lie 及其追随者们（E. Study、G. Scheffers、F. Schur，稍后还有 T. Molien 与 É. Cartan）留下了深刻印象，他们当时正在发展“连续”群理论，尤其是分类问题（特别参见 {[}43{]}，第 387 页）。在 1885--1905 年间，这促使该学派的数学家把与研究李群和李代数时所用的类似方法应用于代数结构的研究。这些方法首先基于考虑代数的元素关于其正则表示的特征多项式（这个多项式已见于上面提到的 Weierstrass 与 Dedekind 的著作中），以及把这个多项式分解为不可约因子。正如 Frobenius 稍后发现的那样，这个分解反映在正则表示分解为不可约分量之中。

在 Lie 学派关于代数的研究过程中，这门理论的“内在”概念逐渐显现。根的概念于 1891 年在 G. Scheffers 的著作 {[}43{]} 中以一种特殊情形（即对根的商是体的直积的情形）出现，并在研究一般情形的 T. Molien {[}31{]} 与 É. Cartan {[}4{]} 的著作中更为清晰（“根”一词来自 Frobenius {[}22{]}）。Study 与 Scheffers {[}43{]} 突出了一个代数是若干个其他代数的直积的概念（Peirce 已有萌芽（{[}38{]}，第 221 页））。最后，Molien {[}31{]} 引入了一个代数的商代数，这一概念本质上等价于双边理想（最先由 É. Cartan 定义 {[}4{]}）或同态（这个名称也来自 Frobenius）。与群论的类比是显然的，后来在 1904 年，Epsteen 与 Wedderburn 考虑了双边理想的合成列并把若尔当--赫尔德定理推广到它们上。这个时期最重要的结果是 T. Molien 的 {[}31{]}：在单群概念的指引下，他定义了（C 上的）单代数并证明这些代数是矩阵代数。接着他证明，C 上任意有限秩代数的结构本质上归结为由 Scheffers 研究过的对根的商是体的直和的情形。此后不久，这些结果被 É. Cartan {[}4{]} 重新发现，并建立得更严格、更清晰。在这个场合，他引入了半单代数的概念，并揭示了与 C 上任意代数相伴的数值不变量，即“É. Cartan 整数”（VIII, 第 180 页，习题 8）。他就这样把这门代数的理论带到了一个此后鲜有进展的地步。\(^{(6)}\) 最后，他把 Molien 的以及他自己关于 R 上代数的结果加以推广。

大约在 1900 年，思想开始朝这样一个方向发展：在与线性代数有关的一切事情中放弃对标量体的任何限制。特别地，我们必须指出以 E. H. Moore 与 L. E. Dickson 为中心的美国学派对有限域研究所给予的强劲推动。这项研究的成果是 Wedderburn 定理 {[}50{]}，它证明每个有限域都是交换的。1907 年，Wedderburn 接过 Cartan 的结果并把它们推广到任意基域上 {[}51{]}。在这样做时，他完全放弃了其前辈们的方法（当体不再是代数闭的或极大序的时候，那些方法不再适用）。他转而对 B. Peirce 的幂等元技巧加以提炼后重新使用，这使他得以把半单代数的结构定理化为确定的形式，把半单代数的研究归结为非交换体的研究。此外，标量扩张的问题在他的视角下自然出现，Wedderburn 证明每个半单代数在基域的可分扩张后仍是半单的 \(^{(7)}\)，且当这个扩张足够大时成为中心单矩阵代数的直积（{[}51{]}，第 102 页）。\(^{(8)}\) 稍后，Dickson 对 n = 3 {[}17{]}，以及 Wedderburn 本人对任意 n {[}52{]}，给出了在其中心上秩为 \(n^{2}\) 的非交换体的最初例子，\(^{(9)}\) 从而在特殊情形引入了后来由 R. Brauer {[}2{]} 与 E. Noether（{[}33{]}、{[}34{]} 与 {[}35{]}）发展的“交叉乘积”与“因子系”理论。最后，在 1921 年，Wedderburn 证明了交换子定理的一个特殊情形 {[}53{]}。

与此同时，从 1896 年到 1910 年，一门与代数理论相近的理论，即群的线性表示理论（起初仅限于有限群的表示），由 Frobenius、Burnside 与 I. Schur 发展起来。它源于 Dedekind 所作的一些评注。在研究伽罗瓦扩张的正规基的过程中，Dedekind（甚至还在他关于代数的著作发表之前）大约在 1880 年遇到了“群行列式” \(\det(x_{st^{-1}})\)，其中 \((x_{s})_{s\in G}\) 是以有限群 G 为指标集的一个变量序列（换句话说，就是群 G 的代数的泛元素关于其正则表示的范数）。他观察到，当 G 交换时，这个多项式分解为线性因子（这推广了人们在很久以前对相应于循环群 G 的“循环”矩阵的行列式所证明的一个恒等式）。在他与 Frobenius 的引人入胜的通信 {[}13{]} 中，Dedekind 于 1896 年使后者注意到这一性质、它与交换群特征标理论的联系，以及他于 1886 年得到的关于某些特定非交换群的类似结果。几个月后，Frobenius 借助他对特征标概念的出色推广 {[}20{]}（我们不在此讨论它）完全解决了“群行列式”分解为不可约因子的问题 {[}19{]}。然而应当指出，在这一理论后来的发展中，\(^{(10)}\) Frobenius 始终意识到它与代数理论的联系（Dedekind 在他的信中从未停止强调这一点）。在引入了群的不可约表示与完全可约表示的概念 {[}21{]} 并证明了正则表示包含所有的不可约表示之后，正是通过类似的方法，他于 1903 年着手重新研究 Molien--Cartan 的理论 {[}22{]}。在 Burnside {[}3{]} 与 I. Schur {[}45{]} 的工作中，这一理论的“超复”方面并未明显地介入；但正是在他们的工作中，出现了不可约表示的基本性质：舒尔引理与伯恩赛德定理。最后还应当指出，正是在这一理论中，交换子定理的两个特殊情形第一次出现。第一个出现在 I. Schur 的学位论文 {[}44{]} 中，该论文（正是通过在一个张量空间的自同态环中的交换）把线性群的表示与对称群的表示联系起来。第二个出现在 Schur 1905 年的工作 {[}46{]} 中，他在其中证明，在域 C 上与一个不可约表示的所有矩阵都交换的矩阵必是 \(I\) 的标量倍（这个结果也可以由伯恩赛德定理推出）。

剩下的工作是把这些理论的共同基础明确地辨认出来：这是 1925--1933 年间以 E. Noether 和 E. Artin 为核心的德国学派所完成的工作，现代代数正是在这一时期创立的。早在 1903 年，H. Poincaré 就在一篇关于线性微分方程的代数积分的论文 {[}42{]} 中，在一个代数里定义了左理想、右理想以及极小理想的概念。他还观察到，在半单代数中，每个左理想都是它与各个单分量的交的直和，并且在 n 阶矩阵代数中，极小理想的维数是 n；但他的工作没有引起代数学家的注意。\(^{(11)}\) 1907 年，Wedderburn 重新定义了代数的左理想与右理想，并证明了它们的若干性质（特别是证明了根是最大的幂零左理想（{[}51{]}，第 113--114 页））。

但直到 1927 年，这些概念才在代数理论中得到本质性的使用。\(^{(12)}\) 通过对此前散见于各处的证明方法加以推广，\(^{(13)}\) W. Krull 在 1925 年 {[}28{]} 与 E. Noether 在 1926 年 {[}33{]} 引入并系统地使用了极大条件与极小条件。前者用它们把关于有限群分解为不可分解群的直积的雷马克定理推广到带算子的交换群（这一概念就是他在此时定义的）（参见 VIII，第 37 页，定理 2），后者则在刻画戴德金环时使用了这些条件。1927 年，E. Artin {[}1{]} 把同一思想应用于非交换环，说明了如何通过对极小理想的系统研究，把韦德伯恩定理推广到其左理想同时满足极大条件与极小条件的所有环。\(^{(14)}\)

在另一方面，1926 年，Krull {[}29{]} 建立了带算子的交换群的概念与群的线性表示概念之间的联系。这一观点被 E. Noether 在 1929 年一篇奠基性的工作 {[}34{]} 中推广到代数并得到详细发挥。鉴于所引入思想的重要性以及阐述的清晰，这篇著作理应与施泰尼茨关于交换域的论文并列为现代代数的支柱之一。\(^{(15)}\)

最后，在始于 1927 年的一系列工作（{[}36{]}、{[}2{]} 与 {[}35{]}）中，E. Noether 与 R. Brauer（1929 年又有 A. Albert 与 H. Hasse 加入其中）在 Wedderburn 与 Dickson 中断之处重新展开对非交换体的研究。虽然他们的结果中最本质的部分在于对布饶尔群、特别是代数数域上的布饶尔群的深入研究，但正是在这些工作中，交换子定理以及单代数的中和域的概念、它与极大交换子域的关系得到了澄清。最后，1927 年，Skolem 刻画了单环的自同构 {[}49{]}，所得的定理几年后又由 E. Noether {[}35{]} 与 R. Brauer {[}2{]} 重新得出。

就这样，到 1934 年，单环与半单环的初等理论已经或多或少达到了它的最终形式（关于这一理论当时状况的综述，见 {[}16{]}）。
% semantic-fragments: legacy-input-seam
\relax{}
